\chapter{Common Frontend Errors}

\begin{notebox}[NOTE]
This chapter is a reference, not a tutorial. Each entry follows the same
pattern --- \textbf{Cause}, \textbf{Identify}, \textbf{Solution} --- so you
can scan straight to the error you are seeing.
\end{notebox}

\subsubsection*{CORS Error}
\noindent\begin{tabularx}{\textwidth}{@{}>{\bfseries}p{2.3cm}X@{}}
\toprule
Cause & Backend response is missing an \texttt{Access-Control-Allow-Origin}
header that matches the frontend's origin. \\
Identify & Console shows "has been blocked by CORS policy"; the Network
tab shows the request was sent but the browser refused the response. \\
Solution & Add \texttt{cors(\{ origin: FRONTEND\_URL, credentials: true \})}
on the backend (Chapter 26). \\
\bottomrule
\end{tabularx}

\subsubsection*{Network Error}
\noindent\begin{tabularx}{\textwidth}{@{}>{\bfseries}p{2.3cm}X@{}}
\toprule
Cause & The request never reached any server --- backend is down, wrong
port, or no internet. \\
Identify & Axios error has no \texttt{response} object, only
\texttt{error.request}; Network tab shows the request as "failed"
with no status. \\
Solution & Confirm the backend is running, check the URL/port, and
verify \texttt{VITE\_API\_URL}. \\
\bottomrule
\end{tabularx}

\subsubsection*{404 on API Call}
\noindent\begin{tabularx}{\textwidth}{@{}>{\bfseries}p{2.3cm}X@{}}
\toprule
Cause & The frontend is calling a path that doesn't exist on the backend
(typo, wrong prefix, or route not mounted). \\
Identify & Network tab shows status \texttt{404} with a response body
(unlike Network Error, the server did respond). \\
Solution & Compare the exact path against \texttt{app.use('/api/...')}
and the router's route definitions. \\
\bottomrule
\end{tabularx}

\subsubsection*{"Failed to fetch"}
\noindent\begin{tabularx}{\textwidth}{@{}>{\bfseries}p{2.3cm}X@{}}
\toprule
Cause & Generic browser message for a request that could not complete ---
server down, CORS preflight failure, or mixed HTTP/HTTPS content. \\
Identify & Thrown by \texttt{fetch}/Axios before any response is received;
check the Network tab for the specific underlying reason (CORS, refused
connection, etc.). \\
Solution & Fix the underlying cause shown in the Network tab --- this
message alone is not diagnostic. \\
\bottomrule
\end{tabularx}

\subsubsection*{"Cannot read properties of undefined (reading 'x')"}
\noindent\begin{tabularx}{\textwidth}{@{}>{\bfseries}p{2.3cm}X@{}}
\toprule
Cause & Rendering \texttt{data.x} before \texttt{data} has arrived from
the API (initial state is \texttt{undefined} or \texttt{null}). \\
Identify & Error points to the exact line accessing the property; happens
on first render, disappears after the fetch resolves. \\
Solution & Initialize state to a sensible default (\texttt{[]}, \texttt{null})
and guard renders with \texttt{data?.x} or an early loading return. \\
\bottomrule
\end{tabularx}

\subsubsection*{State Not Updating (Stale Closures)}
\noindent\begin{tabularx}{\textwidth}{@{}>{\bfseries}p{2.3cm}X@{}}
\toprule
Cause & A callback/effect captured an old value of state from a previous
render and keeps using it. \\
Identify & Logging state inside an event handler shows the value from
several renders ago, not the latest. \\
Solution & Use the updater-function form (\texttt{setCount(c => c + 1)})
or add the value to the effect's dependency array. \\
\bottomrule
\end{tabularx}

\subsubsection*{useEffect Infinite Loop}
\noindent\begin{tabularx}{\textwidth}{@{}>{\bfseries}p{2.3cm}X@{}}
\toprule
Cause & The effect updates state that is also in its own dependency array,
or depends on a new object/array/function created every render. \\
Identify & Network tab shows the same request firing repeatedly;
browser becomes sluggish. \\
Solution & Remove the unstable dependency, wrap it in
\texttt{useMemo}/\texttt{useCallback}, or move the state update outside
the effect's own trigger condition. \\
\bottomrule
\end{tabularx}

\subsubsection*{Duplicate API Calls}
\noindent\begin{tabularx}{\textwidth}{@{}>{\bfseries}p{2.3cm}X@{}}
\toprule
Cause & \texttt{React.StrictMode} intentionally double-invokes effects in
development, or the effect is missing a dependency and re-fires on every
render. \\
Identify & Two identical requests in the Network tab per page load;
persists in dev only (StrictMode) or also in production (real bug). \\
Solution & If it only happens in dev, it's expected --- ignore it. If it
also happens in production, fix the effect's dependency array. \\
\bottomrule
\end{tabularx}

\subsubsection*{React Key Warning}
\noindent\begin{tabularx}{\textwidth}{@{}>{\bfseries}p{2.3cm}X@{}}
\toprule
Cause & A list rendered with \texttt{.map()} is missing a \texttt{key}
prop, or two items share the same key (often \texttt{index} used after a
reorder/filter). \\
Identify & Console warning: "Each child in a list should have a unique
\texttt{key} prop." \\
Solution & Use a stable unique value, typically \texttt{task.\_id}, never
the array index for lists that can reorder. \\
\bottomrule
\end{tabularx}

\subsubsection*{Component Not Rendering}
\noindent\begin{tabularx}{\textwidth}{@{}>{\bfseries}p{2.3cm}X@{}}
\toprule
Cause & A conditional (\texttt{if}, \texttt{\&\&}, early \texttt{return
null}) evaluates falsy when it shouldn't, or the component isn't mounted
in the tree at all. \\
Identify & No error in console --- the component simply never appears;
React DevTools shows it missing from the tree. \\
Solution & Log the condition's value right before the branch, and confirm
the parent actually renders this child. \\
\bottomrule
\end{tabularx}

\subsubsection*{Wrong API URL}
\noindent\begin{tabularx}{\textwidth}{@{}>{\bfseries}p{2.3cm}X@{}}
\toprule
Cause & \texttt{VITE\_API\_URL} is unset (falls back to \texttt{undefined})
or has a trailing slash that doubles up with a leading slash on the
request path. \\
Identify & Request URL in the Network tab shows
\texttt{undefined/api/tasks} or \texttt{//api/tasks}. \\
Solution & Set \texttt{VITE\_API\_URL} in \texttt{.env}, keep it without a
trailing slash, and build paths as \texttt{`\$\{baseURL\}/api/tasks`}. \\
\bottomrule
\end{tabularx}

\subsubsection*{Environment Variable Undefined}
\noindent\begin{tabularx}{\textwidth}{@{}>{\bfseries}p{2.3cm}X@{}}
\toprule
Cause & Vite only exposes env variables prefixed with \texttt{VITE\_} to
client code, or the dev server was started before the \texttt{.env} file
was created/edited. \\
Identify & \texttt{import.meta.env.VITE\_API\_URL} logs \texttt{undefined}
even though the variable exists in \texttt{.env}. \\
Solution & Rename the variable with the \texttt{VITE\_} prefix and restart
\texttt{npm run dev} --- Vite does not hot-reload \texttt{.env} changes. \\
\bottomrule
\end{tabularx}

\begin{tipbox}[TIP]
Every one of these errors is diagnosed the same way: open DevTools,
check the Console tab for the message, then the Network tab for the
actual request/response. Guessing without opening DevTools wastes time.
\end{tipbox}
